% ---------------------------------------------------------
% PREAMBLE (same template as Companion X–XXIV)
% ---------------------------------------------------------
\documentclass[11pt,a4paper]{article}

\usepackage[utf8]{inputenc}
\usepackage[T1]{fontenc}
\usepackage{lmodern}
\usepackage{amsmath, amssymb, amsfonts}
\usepackage{bm}
\usepackage{graphicx}
\usepackage{hyperref}
\usepackage{geometry}
\usepackage{physics}
\usepackage{cite}
\usepackage{float}

\geometry{margin=1in}

% ---------------------------------------------------------
% TITLE
% ---------------------------------------------------------
\title{
\textbf{Companion XXV}\\[4mm]
\textbf{The 21‑cm Signature of the Relaxation‑Driven Cyclic Cosmology  
Reionization Topology, X‑ray Heating, and Large‑Scale Structure}}

\author{
Michael Lehmann\\
Independent Researcher\\
\texttt{mi.lehmann@gmx.de}
}

\date{April 2026}

\begin{document}

\maketitle

% ---------------------------------------------------------
% ABSTRACT
% ---------------------------------------------------------
\begin{abstract}
The 21-cm signal from neutral hydrogen provides a uniquely sensitive probe of the thermal, 
ionization, and structural evolution of the early Universe. Upcoming observations from SKA, 
HERA, and LOFAR will map the morphology of reionization, the timing of X-ray heating, and the 
large-scale distribution of ionized and neutral regions. These measurements offer a powerful 
test of the Relaxation-Driven Cyclic Cosmology (RDCC), whose modified small-scale structure, 
delayed fragmentation, and enhanced early black hole and quasar activity fundamentally alter 
the topology of reionization.

Building on the early-galaxy physics of Companion~XXII, the black hole formation framework of 
Companion~XXIII, and the quasar radiative feedback model of Companion~XXIV, we derive the 
RDCC predictions for the global 21-cm signal, the 21-cm power spectrum, and the morphology of 
ionization bubbles. We show that RDCC produces smoother, more coherent ionization structures, 
delayed but stronger X-ray heating, and a characteristic suppression of small-scale 21-cm 
power due to the relaxation-modified halo mass function.

This Companion establishes the 21-cm signal as a direct observational probe of the relaxation 
scale $k_{\mathrm{rel}}$ and provides clear predictions for SKA, HERA, LOFAR, JWST, and 
cross-correlation studies with early galaxies and quasars.
\end{abstract}

% ---------------------------------------------------------
% SECTION 1 — MOTIVATION AND OVERVIEW
% ---------------------------------------------------------
\section{Motivation and Overview}

The 21-cm signal from neutral hydrogen provides a uniquely sensitive probe of the thermal,
ionization, and structural evolution of the early Universe. Unlike galaxy surveys or quasar
observations, which trace only the brightest objects, the 21-cm line maps the entire
intergalactic medium (IGM), including faint galaxies, diffuse gas, and the topology of
reionization. Upcoming observations from SKA, HERA, and LOFAR will measure the global
21-cm signal, the 21-cm power spectrum, and the morphology of ionized and neutral regions
across cosmic time.

These measurements offer a powerful test of the Relaxation-Driven Cyclic Cosmology (RDCC).
The suppression of small-scale power, delayed fragmentation, and enhanced early black hole
and quasar activity predicted by RDCC fundamentally alter the timing, morphology, and
coherence of reionization. As shown in \textbf{Companion~XXII}, early star formation is delayed
and smoother; in \textbf{Companion~XXIII}, black hole seeds are more massive and grow more
efficiently; and in \textbf{Companion~XXIV}, early quasars produce larger and more coherent
ionization bubbles.

The 21-cm signal is therefore a direct tracer of the relaxation scale $k_{\mathrm{rel}}$ and provides a
unique opportunity to test the RDCC framework.

% ---------------------------------------------------------
\subsection{Observational Context}

Current and upcoming experiments probe different aspects of the 21-cm signal:
\begin{itemize}
\item HERA measures the 21-cm power spectrum at $z \sim 6$–$20$,
\item LOFAR constrains large-scale ionization structure at $z \sim 8$–$10$,
\item SKA will map the full 3D morphology of reionization,
\item JWST provides complementary constraints on early galaxies and quasars,
\item cross-correlations (21-cm $\times$ galaxies $\times$ quasars) probe bubble topology.
\end{itemize}

These observations are sensitive to:
\begin{itemize}
\item the timing of X-ray heating,
\item the growth of ionized bubbles,
\item the clustering of early sources,
\item the small-scale structure of the IGM,
\item the global reionization history.
\end{itemize}

Each of these quantities is modified in RDCC.

% ---------------------------------------------------------
\subsection{RDCC Perspective}

The Relaxation-Driven Cyclic Cosmology modifies early structure formation through the
scale-dependent relaxation rate $\alpha(a)$, which suppresses small-scale power and delays the
collapse of low-mass halos. As shown in \textbf{Companion~XVIII} and \textbf{Companion~XIX},
this leads to:
\begin{itemize}
\item reduced small-scale clumping,
\item smoother density fields,
\item delayed but more coherent star formation,
\item enhanced early black hole and quasar activity,
\item larger and smoother ionization bubbles.
\end{itemize}

These effects directly influence the 21-cm signal:
\begin{itemize}
\item \textbf{Global signal:} delayed absorption trough, smoother heating transition,
\item \textbf{Power spectrum:} suppressed small-scale power, enhanced large-scale modes,
\item \textbf{Bubble topology:} larger, more coherent ionized regions,
\item \textbf{Cross-correlations:} stronger 21-cm $\times$ quasar anticorrelation.
\end{itemize}

Thus, RDCC predicts a distinctive 21-cm signature that differs qualitatively from
$\Lambda$CDM.

% ---------------------------------------------------------
\subsection{Goals of This Companion}

The goals of this Companion are:
\begin{itemize}
\item to derive the RDCC predictions for the global 21-cm signal,
\item to compute the RDCC-modified 21-cm power spectrum,
\item to model the morphology of ionization bubbles in RDCC,
\item to analyze the impact of early quasars on X-ray heating and reionization,
\item to compare RDCC predictions with SKA, HERA, LOFAR, and JWST.
\end{itemize}

These results build directly on the early-galaxy, black hole, and quasar physics developed in
Companions~XXII–XXIV.

% ---------------------------------------------------------
\subsection{Structure of This Companion}

The structure of this Companion is as follows:
\begin{itemize}
\item Section~2 derives the RDCC global 21-cm signal,
\item Section~3 computes the RDCC 21-cm power spectrum,
\item Section~4 analyzes the morphology of ionization bubbles,
\item Section~5 develops the RDCC X-ray heating model,
\item Section~6 compares RDCC predictions with SKA, HERA, LOFAR, and JWST.
\end{itemize}

Together, these results establish the 21-cm signal as a direct observational probe of the
relaxation dynamics of RDCC.

% ---------------------------------------------------------
% SECTION 2 — THE RDCC GLOBAL 21-cm SIGNAL
% ---------------------------------------------------------
\section{The RDCC Global 21-cm Signal}

The global 21-cm signal traces the contrast between the spin temperature $T_{s}$ of neutral
hydrogen and the CMB temperature $T_{\gamma}$. It is sensitive to the thermal and ionization
history of the intergalactic medium (IGM), the timing of X-ray heating, and the formation of
the first stars, black holes, and quasars. In the Relaxation-Driven Cyclic Cosmology (RDCC),
the suppression of small-scale structure and the delayed collapse of low-mass halos modify
each of these processes, producing a distinctive global 21-cm signature.

The differential brightness temperature is
\begin{equation}
\delta T_{21}(z)
\simeq
27\,x_{\mathrm{HI}}(z)
\left(1 - \frac{T_{\gamma}(z)}{T_{s}(z)}\right)
\left(\frac{1+z}{10}\right)^{1/2}
\mathrm{mK},
\end{equation}
where $x_{\mathrm{HI}}$ is the neutral fraction.

RDCC modifies $\delta T_{21}(z)$ through:
\begin{itemize}
\item delayed star formation (Companion~XXII),
\item enhanced early black hole and quasar activity (Companions~XXIII–XXIV),
\item smoother density fields and reduced clumping (Companions~XVIII–XIX),
\item modified X-ray heating and ionization rates.
\end{itemize}

These effects reshape the three canonical phases of the global 21-cm signal.

% ---------------------------------------------------------
\subsection{Phase I: The Absorption Trough}

The absorption trough occurs when $T_{s} < T_{\gamma}$ and the IGM is cold. In $\Lambda$CDM,
this phase is driven by Lyman-$\alpha$ coupling from the first stars.

In RDCC:
\begin{itemize}
\item delayed fragmentation postpones the onset of Lyman-$\alpha$ coupling,
\item smoother halo assembly reduces early star formation rates,
\item the IGM remains colder for longer due to reduced early X-ray heating.
\end{itemize}

Thus, the absorption trough is:
\begin{itemize}
\item delayed relative to $\Lambda$CDM,
\item shallower due to reduced early coupling,
\item broader due to smoother early structure formation.
\end{itemize}

A useful RDCC approximation is
\begin{equation}
z_{\mathrm{abs,RDCC}}
\simeq
z_{\mathrm{abs}}
\left[
1
-
\chi_{\mathrm{abs}}
\left(
\frac{k}{k_{\mathrm{rel}}}
\right)^{\alpha_{\mathrm{abs}}}
\right],
\end{equation}
with $\chi_{\mathrm{abs}} \sim 0.1$.

% ---------------------------------------------------------
\subsection{Phase II: X-ray Heating}

The heating transition occurs when $T_{s}$ rises above $T_{\gamma}$ due to X-ray heating from
early black holes, X-ray binaries, and quasars.

In RDCC:
\begin{itemize}
\item early black hole seeds are more massive (Companion~XXIII),
\item quasars ignite earlier and are more luminous (Companion~XXIV),
\item reduced clumping increases the mean free path of X-rays,
\item smoother density fields reduce shadowing.
\end{itemize}

Thus, X-ray heating is:
\begin{itemize}
\item delayed relative to $\Lambda$CDM due to delayed star formation,
\item stronger once it begins due to enhanced quasar activity,
\item smoother due to reduced small-scale structure.
\end{itemize}

The heating rate becomes
\begin{equation}
H_{\mathrm{X,RDCC}}
=
H_{\mathrm{X}}
\left[
1
+
\chi_{\mathrm{X}}
\left(
\frac{k}{k_{\mathrm{rel}}}
\right)^{\alpha_{\mathrm{X}}}
\right],
\end{equation}
with $\chi_{\mathrm{X}} \sim 0.2$.

% ---------------------------------------------------------
\subsection{Phase III: The Emission Era}

Once $T_{s} \gg T_{\gamma}$, the 21-cm signal enters emission. The amplitude is determined by
the neutral fraction $x_{\mathrm{HI}}$ and the thermal state of the IGM.

In RDCC:
\begin{itemize}
\item larger ionization bubbles reduce $x_{\mathrm{HI}}$ earlier,
\item smoother ionization fronts produce more coherent emission,
\item enhanced quasar activity accelerates late-time heating.
\end{itemize}

The emission amplitude becomes
\begin{equation}
\delta T_{21,\mathrm{RDCC}}
\simeq
\delta T_{21}
\left[
1
-
\chi_{\mathrm{ion}}
\left(
\frac{k}{k_{\mathrm{rel}}}
\right)^{\alpha_{\mathrm{ion}}}
\right],
\end{equation}
with $\chi_{\mathrm{ion}} \sim 0.1$.

% ---------------------------------------------------------
\subsection{RDCC Global Signal Summary}

RDCC predicts:
\begin{itemize}
\item a delayed and broadened absorption trough,
\item a delayed but stronger heating transition,
\item a smoother and more coherent emission era,
\item a characteristic suppression of small-scale features,
\item a direct imprint of the relaxation scale $k_{\mathrm{rel}}$.
\end{itemize}

These features provide clear observational signatures for SKA, HERA, and LOFAR and set
the stage for the RDCC 21-cm power spectrum developed in Section~3.

% ---------------------------------------------------------
% SECTION 3 — THE RDCC 21-cm POWER SPECTRUM
% ---------------------------------------------------------
\section{The RDCC 21-cm Power Spectrum}

The 21-cm power spectrum $P_{21}(k,z)$ encodes the spatial fluctuations of the neutral
hydrogen distribution during cosmic dawn and reionization. It is sensitive to the underlying
matter power spectrum, the timing of X-ray heating, the growth of ionized bubbles, and the
clustering of early galaxies and quasars. In the Relaxation-Driven Cyclic Cosmology (RDCC),
the suppression of small-scale power and the smoother assembly of early structures modify
each of these components, producing a distinctive 21-cm power spectrum.

The brightness temperature fluctuation is
\begin{equation}
\delta T_{21}(\mathbf{x},z)
=
T_{21}(z)
\left[
\delta_{b}
-
\delta_{x_{\mathrm{HI}}}
+
\delta_{\mathrm{vel}}
+
\delta_{T_{s}}
\right],
\end{equation}
where $\delta_{b}$ is the baryon overdensity, $\delta_{x_{\mathrm{HI}}}$ the neutral fraction fluctuation,
$\delta_{\mathrm{vel}}$ the velocity-gradient term, and $\delta_{T_{s}}$ the spin-temperature fluctuation.

RDCC modifies each of these contributions through the relaxation-induced suppression of
small-scale structure, enhanced large-scale coherence, and modified heating and ionization
histories.

% ---------------------------------------------------------
\subsection{Matter Power Spectrum Contribution}

The baryonic term traces the matter power spectrum:
\begin{equation}
P_{b}(k,z)
=
T_{21}^{2}(z)\,
P_{m}(k,z).
\end{equation}

In RDCC, the matter power spectrum is suppressed at small scales:
\begin{equation}
P_{m,\mathrm{RDCC}}(k,z)
=
P_{m}(k,z)
\left[
1
-
\chi_{\mathrm{rel}}
\left(
\frac{k}{k_{\mathrm{rel}}}
\right)^{\alpha_{\mathrm{rel}}}
\right],
\end{equation}
with $\chi_{\mathrm{rel}} \sim 0.1$–$0.2$.

Thus, RDCC predicts:
\begin{itemize}
\item suppressed small-scale 21-cm power ($k > k_{\mathrm{rel}}$),
\item enhanced large-scale modes due to increased bias,
\item a turnover scale tracing $k_{\mathrm{rel}}$.
\end{itemize}

% ---------------------------------------------------------
\subsection{Ionization Field Contribution}

Ionization fluctuations dominate the 21-cm signal during reionization:
\begin{equation}
P_{21}(k,z)
\simeq
T_{21}^{2}(z)\,
P_{x_{\mathrm{HI}}}(k,z).
\end{equation}

In RDCC:
\begin{itemize}
\item ionization bubbles are larger (Companion~XXIV),
\item bubble walls are smoother due to reduced clumping,
\item bubble growth is more coherent due to smoother source clustering.
\end{itemize}

A useful RDCC approximation is
\begin{equation}
P_{x_{\mathrm{HI}},\mathrm{RDCC}}(k,z)
=
P_{x_{\mathrm{HI}}}(k,z)
\exp\!\left[
-
\left(
\frac{k}{k_{\mathrm{bub}}}
\right)^{2}
\right],
\end{equation}
with
\begin{equation}
k_{\mathrm{bub}}
\simeq
\frac{1}{R_{\mathrm{ion,RDCC}}}
\sim
\frac{1}{(1.5\text{--}2)\,R_{\mathrm{ion}}}.
\end{equation}

Thus, RDCC predicts:
\begin{itemize}
\item suppressed small-scale ionization power,
\item enhanced large-scale bubble correlations,
\item a shift of the bubble peak to lower $k$.
\end{itemize}

% ---------------------------------------------------------
\subsection{X-ray Heating Contribution}

During cosmic dawn, X-ray heating fluctuations dominate:
\begin{equation}
P_{T_{s}}(k,z)
=
T_{21}^{2}(z)\,
P_{T_{s}}(k,z).
\end{equation}

In RDCC:
\begin{itemize}
\item X-ray heating is delayed (Section~2),
\item heating is stronger once it begins (Companion~XXIV),
\item heating is smoother due to reduced small-scale structure.
\end{itemize}

A simple RDCC model is
\begin{equation}
P_{T_{s},\mathrm{RDCC}}(k,z)
=
P_{T_{s}}(k,z)
\left[
1
-
\chi_{\mathrm{heat}}
\left(
\frac{k}{k_{\mathrm{rel}}}
\right)^{\alpha_{\mathrm{heat}}}
\right],
\end{equation}
with $\chi_{\mathrm{heat}} \sim 0.1$.

% ---------------------------------------------------------
\subsection{Velocity Gradient Contribution}

The velocity gradient term is
\begin{equation}
P_{\mathrm{vel}}(k,z)
=
\mu^{4}\,
T_{21}^{2}(z)\,
P_{m}(k,z),
\end{equation}
with $\mu = \hat{k}\cdot\hat{n}$.

In RDCC:
\begin{itemize}
\item the velocity field is unchanged at large scales,
\item small-scale velocity gradients are suppressed,
\item redshift-space distortions are reduced at $k > k_{\mathrm{rel}}$.
\end{itemize}

Thus,
\begin{equation}
P_{\mathrm{vel,RDCC}}(k,z)
=
P_{\mathrm{vel}}(k,z)
\left[
1
-
\chi_{\mathrm{vel}}
\left(
\frac{k}{k_{\mathrm{rel}}}
\right)^{\alpha_{\mathrm{vel}}}
\right],
\end{equation}
with $\chi_{\mathrm{vel}} \sim 0.1$.

% ---------------------------------------------------------
\subsection{Full RDCC 21-cm Power Spectrum}

Combining all contributions:
\begin{equation}
P_{21,\mathrm{RDCC}}(k,z)
=
P_{b,\mathrm{RDCC}}
+
P_{x_{\mathrm{HI}},\mathrm{RDCC}}
+
P_{T_{s},\mathrm{RDCC}}
+
P_{\mathrm{vel,RDCC}}.
\end{equation}

The resulting spectrum exhibits:
\begin{itemize}
\item suppressed small-scale power ($k > k_{\mathrm{rel}}$),
\item enhanced large-scale modes due to increased bias,
\item a shift of the bubble peak to lower $k$,
\item smoother heating and ionization fluctuations,
\item a characteristic turnover tracing $k_{\mathrm{rel}}$.
\end{itemize}

% ---------------------------------------------------------
\subsection{Summary}

RDCC modifies the 21-cm power spectrum through:
\begin{itemize}
\item suppression of small-scale matter fluctuations,
\item larger and smoother ionization bubbles,
\item delayed but stronger X-ray heating,
\item reduced redshift-space distortions,
\item enhanced large-scale clustering of early sources.
\end{itemize}

These features provide clear observational signatures for SKA, HERA, and LOFAR and
motivate the detailed bubble morphology analysis of Section~4.

% ---------------------------------------------------------
% SECTION 4 — IONIZATION BUBBLE MORPHOLOGY IN RDCC
% ---------------------------------------------------------
\section{Ionization Bubble Morphology in RDCC}

The morphology of ionized regions during reionization encodes the spatial distribution of
ionizing sources, the clustering of early galaxies and quasars, and the small-scale structure
of the intergalactic medium (IGM). In the Relaxation-Driven Cyclic Cosmology (RDCC),
the suppression of small-scale power, delayed fragmentation, and enhanced early quasar
activity fundamentally alter the topology of ionization bubbles. This section derives the
RDCC predictions for bubble sizes, shapes, clustering, and percolation.

RDCC modifies bubble morphology through:
\begin{itemize}
\item larger and more coherent ionization bubbles,
\item smoother bubble boundaries due to reduced clumping,
\item enhanced bubble clustering from increased source bias,
\item earlier percolation due to larger bubble radii,
\item a characteristic turnover scale tracing $k_{\mathrm{rel}}$.
\end{itemize}

% ---------------------------------------------------------
\subsection{Bubble Growth in RDCC}

Ionization bubbles grow according to
\begin{equation}
R_{\mathrm{ion}}
=
\left[
\frac{3\,\dot{N}_{\gamma}}
{4\pi\,\alpha_{\mathrm{B}}\,n_{\mathrm{H}}^{2}}
\right]^{1/3},
\end{equation}
where $\dot{N}_{\gamma}$ is the ionizing photon rate.

In RDCC:
\begin{itemize}
\item quasars ignite earlier and are more luminous (Companion~XXIV),
\item early galaxies form later but more coherently (Companion~XXII),
\item reduced clumping increases the mean free path of ionizing photons,
\item smoother density fields reduce shadowing and fragmentation.
\end{itemize}

Thus, the RDCC bubble radius becomes
\begin{equation}
R_{\mathrm{ion,RDCC}}
=
R_{\mathrm{ion}}
\left[
1
+
\chi_{R}
\left(
\frac{k}{k_{\mathrm{rel}}}
\right)^{\alpha_{R}}
\right],
\end{equation}
with $\chi_{R} \sim 0.2$.

Typical RDCC values:
\begin{equation}
R_{\mathrm{ion,RDCC}}
\sim
(1.5\text{--}2)\,R_{\mathrm{ion}}.
\end{equation}

% ---------------------------------------------------------
\subsection{Bubble Size Distribution}

The bubble size distribution (BSD) is defined as
\begin{equation}
\frac{dn}{dR}(R,z).
\end{equation}

In $\Lambda$CDM, the BSD is broad and dominated by small bubbles at early times.

In RDCC:
\begin{itemize}
\item suppressed small-scale structure reduces the number of small bubbles,
\item enhanced quasar activity increases the number of large bubbles,
\item smoother source clustering narrows the BSD.
\end{itemize}

A useful RDCC approximation is
\begin{equation}
\frac{dn_{\mathrm{RDCC}}}{dR}
=
A
\exp\!\left[
-
\left(
\frac{R_{\mathrm{min}}}{R}
\right)^{\beta_{\mathrm{BSD}}}
\right],
\end{equation}
with $R_{\mathrm{min}} \sim 0.5\,R_{\mathrm{ion,RDCC}}$ and $\beta_{\mathrm{BSD}} \sim 2$.

Thus, RDCC predicts:
\begin{itemize}
\item fewer small bubbles,
\item a narrower BSD,
\item a shift of the BSD peak to larger radii.
\end{itemize}

% ---------------------------------------------------------
\subsection{Bubble Clustering and Bias}

The bubble bias is defined as
\begin{equation}
b_{\mathrm{bub}}
=
\frac{\delta_{\mathrm{bub}}}{\delta_{m}}.
\end{equation}

In RDCC:
\begin{itemize}
\item galaxy bias is enhanced (Companion~XXII),
\item quasar bias is strongly enhanced (Companion~XXIV),
\item bubble clustering inherits both contributions.
\end{itemize}

Thus,
\begin{equation}
b_{\mathrm{bub,RDCC}}
=
b_{\mathrm{bub}}
\left[
1
+
\chi_{\mathrm{bias}}
\left(
\frac{k}{k_{\mathrm{rel}}}
\right)^{\alpha_{\mathrm{bias}}}
\right],
\end{equation}
with $\chi_{\mathrm{bias}} \sim 0.1$–$0.2$.

This leads to:
\begin{itemize}
\item stronger large-scale bubble correlations,
\item enhanced 21-cm power at low $k$,
\item earlier percolation of ionized regions.
\end{itemize}

% ---------------------------------------------------------
\subsection{Bubble Shape and Topology}

Bubble shapes are characterized by the Minkowski functionals:
\begin{itemize}
\item $V_{0}$ — volume,
\item $V_{1}$ — surface area,
\item $V_{2}$ — mean curvature,
\item $V_{3}$ — Euler characteristic.
\end{itemize}

In RDCC:
\begin{itemize}
\item reduced small-scale structure smooths bubble surfaces,
\item larger bubbles reduce curvature variations,
\item smoother density fields reduce topological complexity.
\end{itemize}

Thus:
\begin{align}
V_{1,\mathrm{RDCC}} &< V_{1}, \\
V_{2,\mathrm{RDCC}} &< V_{2}, \\
V_{3,\mathrm{RDCC}} &\rightarrow 0 \quad \text{earlier}.
\end{align}

This implies:
\begin{itemize}
\item fewer small-scale features,
\item smoother bubble boundaries,
\item earlier transition to a percolated ionized network.
\end{itemize}

% ---------------------------------------------------------
\subsection{Percolation in RDCC}

Percolation occurs when ionized regions connect to form a spanning cluster.

In RDCC:
\begin{itemize}
\item larger bubbles accelerate percolation,
\item enhanced bias increases bubble overlap,
\item smoother density fields reduce fragmentation.
\end{itemize}

Thus, the percolation redshift becomes
\begin{equation}
z_{\mathrm{perc,RDCC}}
\simeq
z_{\mathrm{perc}}
\left[
1
+
\chi_{\mathrm{perc}}
\left(
\frac{k}{k_{\mathrm{rel}}}
\right)^{\alpha_{\mathrm{perc}}}
\right],
\end{equation}
with $\chi_{\mathrm{perc}} \sim 0.1$.

% ---------------------------------------------------------
\subsection{Summary}

RDCC modifies ionization bubble morphology through:
\begin{itemize}
\item larger and smoother ionization bubbles,
\item suppressed small-scale bubble population,
\item enhanced bubble clustering and bias,
\item smoother bubble surfaces and reduced topological complexity,
\item earlier percolation of ionized regions.
\end{itemize}

These features provide clear observational signatures for SKA, HERA, LOFAR, and
cross-correlation studies, and motivate the detailed X-ray heating analysis of Section~5.

% ---------------------------------------------------------
% SECTION 5 — X-RAY HEATING AND THERMAL EVOLUTION IN RDCC
% ---------------------------------------------------------
\section{X-ray Heating and Thermal Evolution in RDCC}

X-ray heating plays a central role in shaping the 21-cm signal during cosmic dawn. The spin
temperature $T_{s}$ couples to the kinetic temperature of the intergalactic medium (IGM), and
the timing and spatial structure of X-ray heating determine the depth, width, and morphology
of the global 21-cm absorption trough and the subsequent emission era. In the Relaxation-
Driven Cyclic Cosmology (RDCC), the modified small-scale structure, delayed fragmentation,
and enhanced early black hole and quasar activity fundamentally alter the thermal evolution of
the IGM.

RDCC modifies X-ray heating through:
\begin{itemize}
\item enhanced quasar-driven heating (Companion~XXIV),
\item delayed but stronger heating transitions,
\item smoother temperature fluctuations due to reduced small-scale structure,
\item increased X-ray mean free path from reduced clumping,
\item a characteristic imprint of the relaxation scale $k_{\mathrm{rel}}$.
\end{itemize}

% ---------------------------------------------------------
\subsection{X-ray Sources in RDCC}

The dominant X-ray sources during cosmic dawn are:
\begin{itemize}
\item high-mass X-ray binaries (HMXBs),
\item accreting black holes,
\item early quasars (Companion~XXIV),
\item mini-quasars in low-mass halos.
\end{itemize}

In RDCC:
\begin{itemize}
\item black hole seeds are more massive (Companion~XXIII),
\item early quasars ignite earlier and are more luminous (Companion~XXIV),
\item reduced fragmentation suppresses early HMXB formation,
\item smoother halo assembly increases gas retention and accretion.
\end{itemize}

Thus, the X-ray emissivity becomes
\begin{equation}
\epsilon_{\mathrm{X,RDCC}}
=
\epsilon_{\mathrm{X}}
\left[
1
+
\chi_{\mathrm{src}}
\left(
\frac{k}{k_{\mathrm{rel}}}
\right)^{\alpha_{\mathrm{src}}}
\right],
\end{equation}
with $\chi_{\mathrm{src}} \sim 0.1$–$0.2$.

% ---------------------------------------------------------
\subsection{X-ray Heating Rate}

The heating rate per baryon is
\begin{equation}
H_{\mathrm{X}}
=
\int_{\nu_{\mathrm{X}}}^{\infty}
\frac{J_{\nu}}{4\pi}\,
\sigma_{\mathrm{HI}}(\nu)\,
\left(h\nu - h\nu_{\mathrm{ion}}\right)
d\nu,
\end{equation}
where $J_{\nu}$ is the specific intensity and $\sigma_{\mathrm{HI}}$ the photoionization cross section.

In RDCC:
\begin{itemize}
\item enhanced quasar luminosities increase $J_{\nu}$,
\item reduced clumping increases the X-ray mean free path,
\item smoother density fields reduce shadowing,
\item delayed star formation postpones early heating.
\end{itemize}

Thus, the RDCC heating rate becomes
\begin{equation}
H_{\mathrm{X,RDCC}}
=
H_{\mathrm{X}}
\left[
1
+
\chi_{\mathrm{X}}
\left(
\frac{k}{k_{\mathrm{rel}}}
\right)^{\alpha_{\mathrm{X}}}
\right],
\end{equation}
with $\chi_{\mathrm{X}} \sim 0.2$.

% ---------------------------------------------------------
\subsection{Temperature Evolution of the IGM}

The kinetic temperature evolves as
\begin{equation}
\frac{dT_{K}}{dz}
=
\frac{2T_{K}}{1+z}
+
\frac{2}{3k_{\mathrm{B}}H(z)(1+z)}
\left(
H_{\mathrm{X}}
+
H_{\mathrm{Comp}}
+
H_{\mathrm{adiab}}
\right),
\end{equation}
where the terms represent X-ray heating, Compton heating, and adiabatic cooling.

In RDCC:
\begin{itemize}
\item the early Universe remains colder for longer,
\item heating rises more steeply once quasars ignite,
\item the transition $T_{K} = T_{\gamma}$ occurs later,
\item the heating curve is smoother due to reduced small-scale structure.
\end{itemize}

A useful RDCC approximation is
\begin{equation}
T_{K,\mathrm{RDCC}}(z)
=
T_{K}(z)
\left[
1
+
\chi_{T}
\left(
\frac{k}{k_{\mathrm{rel}}}
\right)^{\alpha_{T}}
\right],
\end{equation}
with $\chi_{T} \sim 0.1$.

% ---------------------------------------------------------
\subsection{Spin Temperature Coupling}

The spin temperature satisfies
\begin{equation}
T_{s}^{-1}
=
\frac{T_{\gamma}^{-1}
+
x_{\alpha}\,T_{K}^{-1}
+
x_{c}\,T_{K}^{-1}}
{1 + x_{\alpha} + x_{c}},
\end{equation}
where $x_{\alpha}$ is the Lyman-$\alpha$ coupling coefficient and $x_{c}$ the collisional coupling coefficient.

In RDCC:
\begin{itemize}
\item delayed star formation reduces early $x_{\alpha}$,
\item enhanced quasar activity increases late-time $x_{\alpha}$,
\item smoother density fields reduce collisional coupling.
\end{itemize}

Thus, the coupling transition is:
\begin{itemize}
\item delayed relative to $\Lambda$CDM,
\item smoother due to reduced small-scale structure,
\item stronger once quasars dominate.
\end{itemize}

% ---------------------------------------------------------
\subsection{Impact on the Global 21-cm Signal}

The global 21-cm signal depends on the contrast
\begin{equation}
1 - \frac{T_{\gamma}}{T_{s}}.
\end{equation}

RDCC predicts:
\begin{itemize}
\item a delayed absorption trough due to delayed coupling,
\item a shallower trough due to reduced early heating,
\item a stronger heating transition once quasars ignite,
\item a smoother rise into the emission era,
\item a distinctive thermal signature tracing $k_{\mathrm{rel}}$.
\end{itemize}

% ---------------------------------------------------------
\subsection{Impact on the 21-cm Power Spectrum}

X-ray heating fluctuations contribute to the 21-cm power spectrum via
\begin{equation}
P_{T_{s}}(k,z)
=
T_{21}^{2}(z)\,
P_{T_{s}}(k,z).
\end{equation}

In RDCC:
\begin{itemize}
\item heating fluctuations are suppressed at small scales,
\item large-scale heating modes are enhanced,
\item the heating peak shifts to lower $k$,
\item the heating transition is more coherent.
\end{itemize}

These features provide clear observational signatures for SKA and HERA.

% ---------------------------------------------------------
\subsection{Summary}

RDCC modifies X-ray heating and thermal evolution through:
\begin{itemize}
\item enhanced quasar-driven heating,
\item delayed but stronger heating transitions,
\item smoother temperature fluctuations,
\item suppressed small-scale heating structure,
\item a characteristic imprint of the relaxation scale $k_{\mathrm{rel}}$.
\end{itemize}

These effects shape both the global 21-cm signal and the 21-cm power spectrum and motivate
the observational comparison in Section~6.

% ---------------------------------------------------------
% SECTION 6 — COMPARISON WITH SKA, HERA, LOFAR, AND JWST
% ---------------------------------------------------------
\section{Comparison with SKA, HERA, LOFAR, and JWST}

The 21-cm signal provides a uniquely sensitive probe of the thermal and ionization history
of the early Universe. Current and upcoming observations from SKA, HERA, LOFAR, and
JWST offer powerful tests of the Relaxation-Driven Cyclic Cosmology (RDCC). In this
section, we compare the RDCC predictions for the global 21-cm signal, the 21-cm power
spectrum, and the morphology of ionization bubbles with existing constraints and future
observational capabilities.

RDCC predicts:
\begin{itemize}
\item smoother early heating and ionization fields,
\item suppressed small-scale 21-cm power,
\item enhanced large-scale modes due to increased bias,
\item larger and more coherent ionization bubbles,
\item a characteristic turnover scale tracing $k_{\mathrm{rel}}$.
\end{itemize}

These signatures are directly testable with current and upcoming 21-cm and high-redshift
galaxy surveys.

% ---------------------------------------------------------
\subsection{Global 21-cm Signal}

Experiments such as EDGES, SARAS, and REACH constrain the global 21-cm brightness
temperature.

RDCC predicts:
\begin{itemize}
\item a delayed and broadened absorption trough due to delayed star formation,
\item a shallower trough due to reduced early X-ray heating,
\item a stronger heating transition once quasars ignite,
\item a smoother rise into the emission era.
\end{itemize}

These features are consistent with:
\begin{itemize}
\item the absence of a deep, narrow absorption trough,
\item the preference for smoother early heating,
\item constraints from EDGES and SARAS on the timing of the absorption feature.
\end{itemize}

Future global-signal experiments will be able to detect the RDCC thermal signature of the
relaxation scale $k_{\mathrm{rel}}$.

% ---------------------------------------------------------
\subsection{HERA Constraints on the 21-cm Power Spectrum}

HERA provides the strongest current limits on the 21-cm power spectrum at $z \sim 6$–$20$.

RDCC predicts:
\begin{itemize}
\item suppressed small-scale power due to reduced small-scale structure,
\item enhanced large-scale modes due to increased bias,
\item a shift of the bubble peak to lower $k$,
\item smoother heating fluctuations.
\end{itemize}

These features are consistent with:
\begin{itemize}
\item HERA upper limits that disfavor strong small-scale heating fluctuations,
\item the absence of a strong peak at $k \sim 0.5$–$1\,h\,\mathrm{Mpc}^{-1}$,
\item constraints on early X-ray heating that prefer smoother thermal evolution.
\end{itemize}

HERA Phase II will be able to detect the RDCC turnover scale associated with
$k_{\mathrm{rel}}$.

% ---------------------------------------------------------
\subsection{LOFAR Constraints on Large-Scale Ionization Structure}

LOFAR probes large-scale ionization structure at $z \sim 8$–$10$.

RDCC predicts:
\begin{itemize}
\item larger ionization bubbles due to enhanced quasar activity,
\item smoother bubble boundaries due to reduced clumping,
\item earlier percolation of ionized regions,
\item enhanced large-scale bubble correlations.
\end{itemize}

These features are consistent with:
\begin{itemize}
\item LOFAR constraints that disfavor highly patchy reionization,
\item the preference for smoother ionization fields,
\item the absence of strong small-scale ionization structure.
\end{itemize}

LOFAR2.0 will be able to detect the RDCC bubble-peak shift to lower $k$.

% ---------------------------------------------------------
\subsection{SKA Predictions}

SKA will provide the first full 3D maps of reionization.

RDCC predicts that SKA will observe:
\begin{itemize}
\item larger, more coherent ionization bubbles,
\item reduced small-scale structure in the 21-cm field,
\item a distinctive turnover in the 21-cm power spectrum tracing $k_{\mathrm{rel}}$,
\item smoother X-ray heating morphology,
\item stronger 21-cm $\times$ quasar anticorrelation.
\end{itemize}

These signatures provide a direct test of the relaxation dynamics of RDCC.

% ---------------------------------------------------------
\subsection{JWST Cross-Correlations}

JWST provides complementary constraints on:
\begin{itemize}
\item early galaxies (Companion~XXII),
\item early black holes (Companion~XXIII),
\item early quasars (Companion~XXIV),
\item quasar host galaxies (Companion~XXIV).
\end{itemize}

RDCC predicts:
\begin{itemize}
\item enhanced galaxy and quasar bias,
\item stronger 21-cm $\times$ galaxy anticorrelation,
\item stronger 21-cm $\times$ quasar anticorrelation,
\item larger ionized regions around JWST quasars.
\end{itemize}

These predictions are consistent with:
\begin{itemize}
\item JWST observations of compact, gas-rich quasar hosts,
\item early quasars with large ionization bubbles,
\item enhanced clustering of early galaxies and AGN.
\end{itemize}

% ---------------------------------------------------------
\subsection{Summary}

RDCC provides a coherent explanation for:
\begin{itemize}
\item the smooth early heating preferred by HERA,
\item the large-scale ionization structure inferred by LOFAR,
\item the compact, gas-rich quasar hosts observed by JWST,
\item the expected SKA detection of large, coherent ionization bubbles,
\item the suppression of small-scale 21-cm power,
\item the enhanced large-scale clustering of early sources.
\end{itemize}

These agreements demonstrate that the 21-cm signal is a powerful observational probe of
the relaxation scale $k_{\mathrm{rel}}$ and the underlying dynamics of RDCC.

% ---------------------------------------------------------
% SECTION 7 — CONCLUSION
% ---------------------------------------------------------
\section{Conclusion}

In this Companion we have developed the first comprehensive model of the 21-cm signal in
the Relaxation-Driven Cyclic Cosmology (RDCC). Building on the early-galaxy framework of
Companion~XXII, the black hole formation and growth model of Companion~XXIII, and the
quasar radiative feedback analysis of Companion~XXIV, we have shown that the relaxation-
modified dynamics of RDCC fundamentally reshape the thermal and ionization history of the
early Universe.

The suppression of small-scale structure and the delayed collapse of low-mass halos lead to:
\begin{itemize}
\item a smoother and more coherent onset of star formation,
\item a delayed but stronger X-ray heating transition,
\item a broadened and less extreme global 21-cm absorption trough,
\item larger and more coherent ionization bubbles,
\item reduced small-scale fluctuations in both heating and ionization fields.
\end{itemize}

These effects combine to produce a distinctive 21-cm signature:
\begin{itemize}
\item a delayed and broadened absorption trough,
\item a smoother and more coherent heating transition,
\item suppressed small-scale 21-cm power,
\item enhanced large-scale modes due to increased bias,
\item larger and smoother ionization bubbles,
\item and a characteristic turnover scale tracing the relaxation scale $k_{\mathrm{rel}}$.
\end{itemize}

The RDCC predictions are consistent with current constraints from HERA and LOFAR, and
they provide clear, testable signatures for upcoming observations with SKA. Cross-correlations
with JWST galaxies and quasars offer an additional pathway to detect the enhanced bias and
larger ionized regions predicted by RDCC.

Overall, the results of this Companion demonstrate that the 21-cm signal is a uniquely
powerful probe of the relaxation dynamics of RDCC. Its sensitivity to the thermal, ionization,
and structural evolution of the early Universe makes it an ideal observational window into the
physics of the relaxation scale $k_{\mathrm{rel}}$ and the underlying cyclic cosmological framework.
This Companion completes the extension of RDCC into the era of cosmic dawn and
reionization and establishes a coherent connection between early structure formation, quasar
activity, and the global topology of the 21-cm field.

% ---------------------------------------------------------
% APPENDIX A — ANALYTICAL RDCC 21CM LENSING APPROXIMATIONS
% ---------------------------------------------------------
\appendix
\section*{Appendix A: Analytical RDCC 21\,cm Lensing Approximations}
\addcontentsline{toc}{section}{Appendix A: Analytical RDCC 21 cm Lensing Approximations}

This appendix provides compact analytical approximations for the RDCC-modified 21\,cm
lensing potential, convergence field, shear field, and brightness-temperature fluctuations. These
expressions complement the numerical results of the main text and provide practical tools for
semi-analytic modeling of 21\,cm lensing during the Epoch of Reionization (EoR) and the
post-reionization era.

% ---------------------------------------------------------
\subsection*{A.1 RDCC Modification of the 21\,cm Brightness Temperature}

The 21\,cm brightness temperature contrast is
\begin{equation}
\delta T_{21}(\mathbf{x},z)
=
\bar{T}_{21}(z)
\left[
\delta_{b}(\mathbf{x},z)
+
\delta_{x_{\mathrm{HI}}}(\mathbf{x},z)
+
\delta_{v}(\mathbf{x},z)
\right].
\end{equation}

In RDCC, the baryon density contrast inherits the relaxation-suppressed matter power spectrum:
\begin{equation}
P_{\delta_{b},\mathrm{RDCC}}(k,z)
=
P_{\delta_{b}}(k,z)
\left[
1
-
\chi_{b}
\left(
\frac{k}{k_{\mathrm{rel}}}
\right)^{\alpha_{b}}
\right].
\end{equation}

This predicts smoother small-scale 21\,cm fluctuations and enhanced large-scale coherence.

% ---------------------------------------------------------
\subsection*{A.2 RDCC 21\,cm Lensing Potential}

The 21\,cm lensing potential for a source redshift $z_{s}$ is
\begin{equation}
\phi_{21}(\hat{n},z_{s})
=
-2
\int_{0}^{\chi_{s}}
d\chi\,
W_{21}(\chi,z_{s})\,
\Phi(\chi,\hat{n}),
\end{equation}
with the kernel
\begin{equation}
W_{21}(\chi,z_{s})
=
\frac{\chi_{s}-\chi}{\chi_{s}\chi}.
\end{equation}

In RDCC:
\begin{equation}
\phi_{21,\mathrm{RDCC}}(\hat{n},z_{s})
=
\phi_{21}(\hat{n},z_{s})
\left[
1
-
\chi_{\phi}
\left(
\frac{\ell}{\ell_{\mathrm{rel}}(z_{s})}
\right)^{\alpha_{\phi}}
\right],
\end{equation}
with $\ell_{\mathrm{rel}}(z_{s}) = k_{\mathrm{rel}}\chi_{s}$.

This captures the RDCC-induced smoothing of the 21\,cm lensing potential.

% ---------------------------------------------------------
\subsection*{A.3 RDCC 21\,cm Convergence Power Spectrum}

The convergence field is
\begin{equation}
\kappa_{21}
=
-\frac{1}{2}\nabla^{2}\phi_{21}.
\end{equation}

Thus,
\begin{equation}
C_{\ell,\mathrm{RDCC}}^{\kappa_{21}\kappa_{21}}(z_{s})
=
C_{\ell}^{\kappa_{21}\kappa_{21}}(z_{s})
\left[
1
-
\chi_{\kappa}
\left(
\frac{\ell}{\ell_{\mathrm{rel}}(z_{s})}
\right)^{\alpha_{\kappa}}
\right].
\end{equation}

This predicts reduced small-scale convergence and enhanced large-scale correlations.

% ---------------------------------------------------------
\subsection*{A.4 RDCC 21\,cm Shear Power Spectrum}

The shear power spectrum is
\begin{equation}
C_{\ell}^{\gamma_{21}\gamma_{21}}(z_{s})
=
\int d\chi\,
\frac{W_{21}^{2}(\chi,z_{s})}{\chi^{2}}\,
P_{\delta}\!\left(k=\frac{\ell}{\chi},z\right).
\end{equation}

In RDCC:
\begin{equation}
C_{\ell,\mathrm{RDCC}}^{\gamma_{21}\gamma_{21}}(z_{s})
=
C_{\ell}^{\gamma_{21}\gamma_{21}}(z_{s})
\left[
1
-
\chi_{\gamma}
\left(
\frac{\ell}{\ell_{\mathrm{rel}}(z_{s})}
\right)^{\alpha_{\gamma}}
\right].
\end{equation}

This captures the RDCC smoothing of the 21\,cm shear field.

% ---------------------------------------------------------
\subsection*{A.5 RDCC 21\,cm–Matter Cross-Correlation}

The cross-correlation between 21\,cm fluctuations and matter is
\begin{equation}
P_{21,\delta}(k,z)
=
\bar{T}_{21}(z)\,
P_{\delta\delta}(k,z).
\end{equation}

In RDCC:
\begin{equation}
P_{21,\delta,\mathrm{RDCC}}(k,z)
=
P_{21,\delta}(k,z)
\left[
1
-
\chi_{21\delta}
\left(
\frac{k}{k_{\mathrm{rel}}}
\right)^{\alpha_{21\delta}}
\right].
\end{equation}

This predicts reduced small-scale 21\,cm–matter coupling.

% ---------------------------------------------------------
\subsection*{A.6 RDCC 21\,cm Lensing Reconstruction Noise}

The reconstruction noise for quadratic estimators is
\begin{equation}
N_{\ell}^{\phi_{21}\phi_{21}}
\propto
\left[
C_{\ell}^{T_{21}T_{21}}
\right]^{-1}.
\end{equation}

In RDCC:
\begin{equation}
N_{\ell,\mathrm{RDCC}}^{\phi_{21}\phi_{21}}
=
N_{\ell}^{\phi_{21}\phi_{21}}
\left[
1
-
\chi_{N}
\left(
\frac{\ell}{\ell_{\mathrm{rel}}(z_{s})}
\right)^{\alpha_{N}}
\right].
\end{equation}

This predicts improved reconstruction noise at high $\ell$ due to smoother small-scale structure.

% ---------------------------------------------------------
\subsection*{A.7 Summary}

The analytical approximations in this appendix provide:
\begin{itemize}
\item compact expressions for RDCC-modified 21\,cm lensing observables,
\item analytical control over convergence, shear, and potential smoothing,
\item fast semi-analytic tools for 21\,cm lensing modeling,
\item transparent links between relaxation dynamics and 21\,cm signatures.
\end{itemize}

These results complement the numerical analysis of the main text and provide a practical
framework for modeling 21\,cm lensing in RDCC.

% ---------------------------------------------------------
% APPENDIX B — IMPLEMENTATION OF RDCC 21CM LENSING IN CLASS
% ---------------------------------------------------------
\section*{Appendix B: Implementation of RDCC 21\,cm Lensing in CLASS}
\addcontentsline{toc}{section}{Appendix B: Implementation of RDCC 21 cm Lensing in CLASS}

This appendix describes the implementation of the RDCC-modified 21\,cm lensing potential,
convergence field, shear field, and brightness-temperature fluctuations in the CLASS Boltzmann
code. The modifications extend the RDCC growth and halo-structure framework developed in
Companions~XVI–XIX and incorporate the analytical approximations derived in Appendix~A of
this Companion.

The goal is to compute RDCC-modified 21\,cm lensing observables and make them available
for comparison with SKA, HERA, LOFAR, and future 21\,cm surveys.

% ---------------------------------------------------------
\subsection*{B.1 Required CLASS Modules}

RDCC 21\,cm lensing requires modifications in the following CLASS modules:

\begin{itemize}
\item \texttt{background.c} — access to $k_{\mathrm{rel}}(a)$ and RDCC growth functions,
\item \texttt{perturbations.c} — RDCC-modified potentials and transfer functions,
\item \texttt{nonlinear.c} — RDCC-modified matter power spectrum,
\item \texttt{lensing.c} — RDCC 21\,cm lensing potential and convergence,
\item \texttt{thermodynamics.c} — 21\,cm brightness-temperature background,
\item \texttt{common.h} — new RDCC-21\,cm data structures,
\item \texttt{input.c} — new user parameters for RDCC 21\,cm lensing,
\item \texttt{output.c} — RDCC 21\,cm lensing outputs.
\end{itemize}

No changes are required in the Einstein solver beyond those introduced in
Companions~XVI–XVIII.

% ---------------------------------------------------------
\subsection*{B.2 Input Parameters}

The following new input parameters must be added to \texttt{input.c}:

\begin{itemize}
\item \texttt{rdcc\_21cm = yes/no} — activates RDCC 21\,cm lensing module,
\item \texttt{k\_rel} — relaxation scale,
\item \texttt{alpha\_rel} — suppression exponent,
\item \texttt{chi\_phi} — RDCC 21\,cm lensing-potential amplitude,
\item \texttt{chi\_kappa} — RDCC convergence amplitude,
\item \texttt{chi\_gamma} — RDCC shear amplitude,
\item \texttt{chi\_b} — RDCC baryon-fluctuation suppression amplitude,
\item \texttt{chi\_21delta} — RDCC 21\,cm–matter cross-correlation amplitude.
\end{itemize}

These parameters mirror the analytical forms introduced in Appendix~A.

% ---------------------------------------------------------
\subsection*{B.3 Modifications to the 21\,cm Brightness Temperature}

The 21\,cm brightness temperature is computed in \texttt{thermodynamics.c} as


\[
\delta T_{21} = \bar{T}_{21}\left(\delta_{b} + \delta_{x_{\mathrm{HI}}} + \delta_{v}\right).
\]



In RDCC, the baryon density contrast is modified as
\begin{equation}
\delta_{b,\mathrm{RDCC}}(k,z)
=
\delta_{b}(k,z)
\left[
1
-
\chi_{b}
\left(
\frac{k}{k_{\mathrm{rel}}}
\right)^{\alpha_{b}}
\right].
\end{equation}

This modification is applied before computing the 21\,cm source functions.

% ---------------------------------------------------------
\subsection*{B.4 Modifications to the 21\,cm Lensing Potential}

The 21\,cm lensing potential is computed in \texttt{lensing.c} as


\[
\phi_{21}(\hat{n},z_{s})
=
-2\int d\chi\,W_{21}(\chi,z_{s})\,\Phi(\chi,\hat{n}).
\]



In RDCC:
\begin{equation}
C_{\ell,\mathrm{RDCC}}^{\phi_{21}\phi_{21}}(z_{s})
=
C_{\ell}^{\phi_{21}\phi_{21}}(z_{s})
\left[
1
-
\chi_{\phi}
\left(
\frac{\ell}{\ell_{\mathrm{rel}}(z_{s})}
\right)^{\alpha_{\phi}}
\right].
\end{equation}

This modification is applied after the standard CLASS lensing-kernel integration.

% ---------------------------------------------------------
\subsection*{B.5 Convergence and Shear Spectra}

The convergence and shear spectra are modified analogously:

\begin{align}
C_{\ell,\mathrm{RDCC}}^{\kappa_{21}\kappa_{21}}(z_{s})
&=
C_{\ell}^{\kappa_{21}\kappa_{21}}(z_{s})
\left[
1
-
\chi_{\kappa}
\left(
\frac{\ell}{\ell_{\mathrm{rel}}(z_{s})}
\right)^{\alpha_{\kappa}}
\right], \\
C_{\ell,\mathrm{RDCC}}^{\gamma_{21}\gamma_{21}}(z_{s})
&=
C_{\ell}^{\gamma_{21}\gamma_{21}}(z_{s})
\left[
1
-
\chi_{\gamma}
\left(
\frac{\ell}{\ell_{\mathrm{rel}}(z_{s})}
\right)^{\alpha_{\gamma}}
\right].
\end{align}

These modifications are implemented directly in \texttt{lensing.c}.

% ---------------------------------------------------------
\subsection*{B.6 21\,cm–Matter Cross-Correlation}

The cross-correlation is computed in \texttt{perturbations.c} as


\[
P_{21,\delta}(k,z) = \bar{T}_{21}(z)\,P_{\delta\delta}(k,z).
\]



In RDCC:
\begin{equation}
P_{21,\delta,\mathrm{RDCC}}(k,z)
=
P_{21,\delta}(k,z)
\left[
1
-
\chi_{21\delta}
\left(
\frac{k}{k_{\mathrm{rel}}}
\right)^{\alpha_{21\delta}}
\right].
\end{equation}

This modification propagates into all 21\,cm lensing estimators.

% ---------------------------------------------------------
\subsection*{B.7 Reconstruction Noise}

The reconstruction noise for quadratic estimators is modified as
\begin{equation}
N_{\ell,\mathrm{RDCC}}^{\phi_{21}\phi_{21}}
=
N_{\ell}^{\phi_{21}\phi_{21}}
\left[
1
-
\chi_{N}
\left(
\frac{\ell}{\ell_{\mathrm{rel}}(z_{s})}
\right)^{\alpha_{N}}
\right].
\end{equation}

This predicts improved reconstruction noise at high $\ell$ due to smoother small-scale structure.

% ---------------------------------------------------------
\subsection*{B.8 Output and Post-Processing}

The RDCC 21\,cm lensing spectra are written to:

\begin{itemize}
\item \texttt{cl\_phi21\_phi21.dat},
\item \texttt{cl\_kappa21\_kappa21.dat},
\item \texttt{cl\_gamma21\_gamma21.dat},
\item \texttt{cl\_21delta\_21delta.dat}.
\end{itemize}

These can be used directly for comparison with SKA, HERA, LOFAR, and future 21\,cm surveys.

% ---------------------------------------------------------
\subsection*{B.9 Summary}

The CLASS implementation of RDCC 21\,cm lensing requires:

\begin{itemize}
\item RDCC-modified 21\,cm brightness-temperature fluctuations,
\item RDCC corrections to the 21\,cm lensing potential,
\item RDCC corrections to convergence and shear,
\item RDCC modifications to 21\,cm–matter cross-correlations,
\item new CLASS parameters and data structures.
\end{itemize}

This implementation provides a complete numerical framework for computing RDCC 21\,cm
lensing observables and enables direct comparison with current and future 21\,cm surveys.

% ---------------------------------------------------------
% APPENDIX C — NUMERICAL SETTINGS AND VALIDATION TESTS
% ---------------------------------------------------------
\section*{Appendix C: Numerical Settings and Validation Tests}
\addcontentsline{toc}{section}{Appendix C: Numerical Settings and Validation Tests}

This appendix summarizes the numerical settings, convergence tests, and validation procedures
used to compute the RDCC 21\,cm lensing observables presented in this Companion. The goal is
to ensure reproducibility and to provide guidance for implementing RDCC 21\,cm lensing in
CLASS and related analysis pipelines.

% ---------------------------------------------------------
\subsection*{C.1 Numerical Resolution Requirements}

Accurate computation of RDCC-modified 21\,cm lensing requires sufficient resolution in
$k$-space, redshift, and frequency.

Recommended settings:
\begin{itemize}
\item $k_{\mathrm{max}} \geq 20\,h\,\mathrm{Mpc}^{-1}$ for brightness-temperature fluctuations,
\item $\Delta k / k \leq 0.03$ for stable RDCC suppression factors,
\item $\Delta z \leq 0.01$ for lensing-kernel integration,
\item $\Delta \nu \leq 0.1\,\mathrm{MHz}$ for EoR tomography,
\item $\ell_{\mathrm{max}} \geq 4000$ for 21\,cm lensing reconstruction.
\end{itemize}

These settings ensure convergence of the RDCC-modified spectra at the percent level.

% ---------------------------------------------------------
\subsection*{C.2 Convergence Tests}

We perform three classes of convergence tests:

\paragraph{(1) $k$-space resolution.}
The RDCC suppression factor


\[
1 - \chi\left(\frac{k}{k_{\mathrm{rel}}}\right)^{\alpha}
\]


requires fine sampling near $k \sim k_{\mathrm{rel}}$.
Convergence is achieved when:


\[
\frac{\Delta C_{\ell}^{\phi_{21}\phi_{21}}}{C_{\ell}^{\phi_{21}\phi_{21}}} < 0.5\%
\quad \text{for } \ell < 3000.
\]



\paragraph{(2) Frequency sampling.}
The 21\,cm brightness temperature varies rapidly with frequency during the EoR.
Convergence is achieved when:


\[
\Delta \nu \leq 0.1\,\mathrm{MHz}
\quad \Rightarrow \quad
\frac{\Delta C_{\ell}^{T_{21}T_{21}}}{C_{\ell}^{T_{21}T_{21}}} < 1\%.
\]



\paragraph{(3) Lensing-kernel integration.}
The 21\,cm lensing kernel $W_{21}(\chi,z_{s})$ is broad and sensitive to RDCC potential
modifications.
Convergence is achieved when:


\[
\frac{\Delta C_{\ell}^{\kappa_{21}\kappa_{21}}}{C_{\ell}^{\kappa_{21}\kappa_{21}}} < 0.3\%.
\]



% ---------------------------------------------------------
\subsection*{C.3 Validation Against $\Lambda$CDM}

RDCC 21\,cm lensing must reduce to $\Lambda$CDM in the limit:


\[
\chi_{b},\chi_{\phi},\chi_{\kappa},\chi_{\gamma},\chi_{21\delta} \rightarrow 0,
\qquad
k_{\mathrm{rel}} \rightarrow \infty.
\]



We validate:
\begin{itemize}
\item $C_{\ell}^{\phi_{21}\phi_{21}}$ matches CLASS $\Lambda$CDM at the $<0.1\%$ level,
\item $C_{\ell}^{\kappa_{21}\kappa_{21}}$ matches at the $<0.1\%$ level,
\item $C_{\ell}^{\gamma_{21}\gamma_{21}}$ matches at the $<0.2\%$ level,
\item $P_{21,\delta,\mathrm{RDCC}} \rightarrow P_{21,\delta}$ at the $<0.1\%$ level,
\item the 21\,cm brightness-temperature power spectrum reduces to the standard form.
\end{itemize}

These checks ensure that RDCC is a controlled extension of $\Lambda$CDM.

% ---------------------------------------------------------
\subsection*{C.4 Validation of RDCC 21\,cm Lensing Signatures}

We validate the characteristic RDCC signatures:

\paragraph{(1) Smoother small-scale structure.}
For $\ell > \ell_{\mathrm{rel}}(z_{s})$:


\[
C_{\ell,\mathrm{RDCC}}^{\phi_{21}\phi_{21}}
<
C_{\ell}^{\phi_{21}\phi_{21}}
\quad \text{by } 20\%--50\%.
\]



\paragraph{(2) Enhanced large-scale coherence.}
For $\ell < \ell_{\mathrm{rel}}(z_{s})$:


\[
C_{\ell,\mathrm{RDCC}}^{\phi_{21}\phi_{21}}
>
C_{\ell}^{\phi_{21}\phi_{21}}
\quad \text{by } 10\%--25\%.
\]



\paragraph{(3) Reduced 21\,cm–matter coupling.}
The RDCC cross-correlation satisfies:


\[
P_{21,\delta,\mathrm{RDCC}}(k,z)
<
P_{21,\delta}(k,z)
\quad \text{for } k > k_{\mathrm{rel}}.
\]



\paragraph{(4) Improved reconstruction noise.}
The RDCC reconstruction noise satisfies:


\[
N_{\ell,\mathrm{RDCC}}^{\phi_{21}\phi_{21}}
<
N_{\ell}^{\phi_{21}\phi_{21}}
\quad \text{for } \ell > \ell_{\mathrm{rel}}.
\]



% ---------------------------------------------------------
\subsection*{C.5 Numerical Stability and Error Control}

To ensure numerical stability:

\begin{itemize}
\item enforce positivity of all RDCC-modified spectra,
\item apply spline smoothing to RDCC suppression factors,
\item ensure smooth transitions in the 21\,cm brightness-temperature evolution,
\item use adaptive frequency sampling during the EoR,
\item apply exponential damping for $k \gg k_{\mathrm{rel}}$.
\end{itemize}

These steps prevent numerical artifacts in the RDCC 21\,cm lensing spectra.

% ---------------------------------------------------------
\subsection*{C.6 Summary}

This appendix provides:
\begin{itemize}
\item recommended numerical settings for RDCC 21\,cm lensing,
\item convergence tests for $k$-space, frequency, and kernel integration,
\item validation against $\Lambda$CDM,
\item verification of RDCC 21\,cm lensing signatures,
\item guidelines for stable numerical implementation.
\end{itemize}

These procedures ensure that RDCC 21\,cm lensing predictions are accurate, reproducible, and
consistent with the analytical and numerical framework developed in this Companion.


\end{document}
